% from ega2-1-fr.tex
\oldpage[II]{91}

$\mathcal{J}_i\otimes(\mathcal{L}^{\otimes n}|f^{-1}(U_i))$ est engendré
par ses sections (compte tenu de (\hyperref[I.9.2.2-fr]{I, 9.2.2}) et de
(\hyperref[II.3.4.7-fr]{3.4.7}))~; il suffit donc d'appliquer le critère
(\hyperref[II.4.5.5-fr]{4.5.5, d'')}).

\begin{proposition}[4.6.9]
\label{II.4.6.9-fr}
Soient $f:X\to Y$ un morphisme quasi-compact, $\mathcal{L}$ un
$\mathcal{O}_X$-Module inversible.
\begin{enumerate}
  \item[\rm{(i)}] Soit $n$ un entier $>0$. Pour que $\mathcal{L}$ soit
  $f$-ample, il faut et il suffit que $\mathcal{L}^{\otimes n}$ soit
  $f$-ample.
  \item[\rm{(ii)}] Soit $\mathcal{L}'$ un $\mathcal{O}_X$-Module
  inversible, et supposons qu'il existe un entier $n>0$ tel que
  l'homomorphisme canonique
  $\sigma:f^*(f_*(\mathcal{L}'^{\otimes n}))\to
  \mathcal{L}'^{\otimes n}$ soit surjectif. Alors, si $\mathcal{L}$ est
  $f$-ample, il en est de même de $\mathcal{L}\otimes\mathcal{L}'$.
\end{enumerate}
\end{proposition}

On se ramène en effet aussitôt au cas où $Y$ est affine, et la proposition
est alors conséquence immédiate de (\hyperref[II.4.5.6-fr]{4.5.6}).

\begin{corollary}[4.6.10]
\label{II.4.6.10-fr}
Le produit tensoriel de deux $\mathcal{O}_X$-Modules $f$-amples est
$f$-ample.
\end{corollary}

\begin{proposition}[4.6.11]
\label{II.4.6.11-fr}
Soient $Y$ un préschéma quasi-compact, $f:X\to Y$ un morphisme de type
fini, $\mathcal{L}$ un $\mathcal{O}_X$-Module inversible. Pour que
$\mathcal{L}$ soit $f$-ample, il faut et il suffit qu'il possède une des
propriétés équivalentes suivantes~:
\begin{enumerate}
  \item[\rm{d)}] Il existe $n_0>0$ tel que, pour tout entier $n\geq n_0$,
  $\mathcal{L}^{\otimes n}$ soit très ample relativement à $f$.
  \item[\rm{d')}] Il existe $n>0$ tel que $\mathcal{L}^{\otimes n}$ soit
  très ample relativement à $f$.
\end{enumerate}
\end{proposition}

Si $\mathcal{L}$ est ample relativement à $f$, il y a un recouvrement fini
$(U_i)$ de $Y$ par des ouverts affines tel que les
$\mathcal{L}|f^{-1}(U_i)$ soient amples. On en conclut
(\hyperref[II.4.5.10-fr]{4.5.10}) qu'il existe un entier $n_0$ tel que
$\mathcal{L}^{\otimes n}|f^{-1}(U_i)$ soit très ample relativement à
$f^{-1}(U_i)\to U_i$ pour tout $n\geq n_0$ et tout $i$, donc
$\mathcal{L}^{\otimes n}$ est très ample relativement à $f$
(\hyperref[II.4.4.5-fr]{4.4.5}). Inversement, d') entraîne déjà que
$\mathcal{L}^{\otimes n}$ est $f$-ample
(\hyperref[II.4.6.2-fr]{4.6.2}), donc il en est de même de $\mathcal{L}$
(\hyperref[II.4.6.9-fr]{4.6.9, (i)}).

\begin{corollary}[4.6.12]
\label{II.4.6.12-fr}
Soient $Y$ un préschéma quasi-compact, $f:X\to Y$ un morphisme de type
fini, $\mathcal{L}$, $\mathcal{L}'$ deux $\mathcal{O}_X$-Modules
inversibles. Si $\mathcal{L}$ est $f$-ample, il existe $n_0$ tel que
$\mathcal{L}^{\otimes n}\otimes\mathcal{L}'$ soit très ample relativement
à $f$ pour tout $n\geq n_0$.
\end{corollary}

On raisonne comme dans (\hyperref[II.4.6.11-fr]{4.6.11}) en utilisant un
recouvrement ouvert affine fini de $Y$ et
(\hyperref[II.4.5.11-fr]{4.5.11}).

\begin{proposition}[4.6.13]
\label{II.4.6.13-fr}
\begin{enumerate}
  \item[\rm{(i)}] Pour tout préschéma $Y$, tout
  $\mathcal{O}_Y$-Module inversible $\mathcal{L}$ est ample relativement
  au morphisme identique $1_Y$.
  \item[\rm{(i bis)}] Soient $f:X\to Y$ un morphisme quasi-compact,
  $j:X'\to X$ un morphisme quasi-compact qui est un homéomorphisme de
  l'espace sous-jacent à $X'$ sur un sous-espace de $X$. Si $\mathcal{L}$
  est un $\mathcal{O}_X$-Module ample relativement à $f$,
  $j^*(\mathcal{L})$ est ample relativement à $f\circ j$.
  \item[\rm{(ii)}] Soient $Z$ un préschéma quasi-compact, $f:X\to Y$,
  $g:Y\to Z$ deux morphismes quasi-compacts, $\mathcal{L}$ un
  $\mathcal{O}_X$-Module ample relativement à $f$, $\mathcal{K}$ un
  $\mathcal{O}_Y$-Module ample relativement à $g$. Alors il existe un
  entier $n_0>0$ tel que
  $\mathcal{L}\otimes f^*(\mathcal{K}^{\otimes n})$ soit ample relativement
  à $g\circ f$ pour tout $n\geq n_0$.
  \item[\rm{(iii)}] Soient $f:X\to Y$ un morphisme quasi-compact,
  $g:Y'\to Y$ un morphisme, et posons $X'=X_{(Y')}$. Si $\mathcal{L}$ est
  un $\mathcal{O}_X$-Module ample relativement à $f$,
  $\mathcal{L}'=\mathcal{L}\otimes_Y\mathcal{O}_{Y'}$ est un
  $\mathcal{O}_{X'}$-Module ample relativement à $f_{(Y')}$.
  \item[\rm{(iv)}] Soient $f_i:X_i\to Y_i$ ($i=1,2$) deux
  $S$-morphismes quasi-compacts. Si $\mathcal{L}_i$ est un
  $\mathcal{O}_{X_i}$-Module ample relativement à $f_i$ ($i=1,2$),
  $\mathcal{L}_1\otimes_S\mathcal{L}_2$ est ample relativement à
  $f_1\times_S f_2$.
  \item[\rm{(v)}] Soient $f:X\to Y$, $g:Y\to Z$ deux morphismes tels que
  $g\circ f$ soit quasi-compact. Si un $\mathcal{O}_X$-Module
  $\mathcal{L}$ est ample relativement à $g\circ f$, et si $g$ est séparé
  ou l'espace sous-jacent à $X$ localement noethérien, alors $\mathcal{L}$
  est ample relativement à $f$.

  